\chapter{Known Cases and Failures}\label{known-cases}
\chapterstatus{complete}

\section{Introduction}\label{known-cases:section-introduction}

This chapter collects the cases in which the Hodge conjecture is known, and explains the mechanism in each. They fall into three groups. First, cases where every Hodge class
is forced to be a combination of known algebraic classes: small dimension, cellular varieties, varieties whose Hodge classes are generated by divisors, and hypersurfaces outside the middle
degree. Second, abelian varieties, where the Mumford--Tate group (Section~\ref{hodge-structures:section-mumford-tate}) makes the Hodge classes computable; there the difficulty concentrates in the \emph{Weil classes}.
Third, geometric methods: the decomposition of the diagonal for varieties with small Chow groups, normal functions for cubic fourfolds, and the Kuga--Satake construction for K3 surfaces.
We state precisely what is proved, with proofs of the elementary parts and references for the rest. References: the source \emph{The Hodge Conjecture} (Chapter 6),
\cite{lewis-survey}, \cite{birkenhake-lange}, \cite{voisin-hodge-2}.

\noindent\textbf{Prerequisites.} Section~\ref{hodge-conjecture:section-lefschetz-11} (The Lefschetz (1,1) Theorem), Section~\ref{hodge-structures:section-mumford-tate} (Mumford--Tate Groups), Chapter~\ref{vhs} (Variations of Hodge Structure) and
Section~\ref{vhs:section-intermediate-jacobians} (Intermediate Jacobians and Normal Functions).

Throughout, $X$ is a nonsingular complex projective variety of dimension $n$.

\section{Cases from general theory}\label{known-cases:section-general}

\begin{definition}\label{known-cases:definition-lefschetz-classes}
The \emph{Lefschetz ring} of $X$ is the $\QQ$-subalgebra $\mathrm{Lef}^\bullet(X) \subseteq \bigoplus_pH^{2p}(X; \QQ)$ generated by $\mathrm{Hdg}^1(X)$, the rational classes of divisors. Its elements are
\emph{Lefschetz classes}.
\end{definition}

\begin{proposition}\label{known-cases:proposition-lefschetz-classes}
Lefschetz classes are algebraic. Hence $\mathrm{HC}(X)$ holds whenever $\mathrm{Hdg}^\bullet(X) = \mathrm{Lef}^\bullet(X)$.
\end{proposition}

\begin{proof}
Divisor classes are algebraic by the Lefschetz $(1,1)$-theorem (Theorem~\ref{hodge-conjecture:theorem-lefschetz-11}), and products of algebraic classes are algebraic
(Theorem~\ref{intersection-theory:theorem-cycle-class}(1)).
\end{proof}

\begin{theorem}\label{known-cases:theorem-general}
The Hodge conjecture holds in the following cases:
\begin{enumerate}
\item $p \in \{0, 1, n - 1, n\}$, and hence for $\dim X \leq 3$ (Corollary~\ref{hodge-conjecture:corollary-codimension-n-1});
\item cellular varieties, such as flag varieties and toric varieties (Example~\ref{hodge-conjecture:example-cellular});
\item very general abelian varieties and all their powers (Theorem~\ref{hodge-structures:theorem-very-general});
\item smooth complete intersections, in all degrees except the middle one (Proposition~\ref{known-cases:proposition-hypersurfaces});
\item uniruled fourfolds \cite{conte-murre-1978} and, more generally, varieties whose group of $0$-cycles is supported on a subvariety of dimension $\leq 3$, in codimension $2$
\cite{bloch-srinivas-1983};
\item cubic fourfolds \cite{zucker-1977};
\item Fermat hypersurfaces of many degrees and dimensions \cite{shioda-1979}.
\end{enumerate}
\end{theorem}

Items (1)--(4) are proved in this book. For (5)--(7) see the references and Sections~\ref{known-cases:section-diagonal} and~\ref{known-cases:section-hypersurfaces}.

\begin{proposition}\label{known-cases:proposition-cellular-factor}
Let $X$ and $Y$ be nonsingular projective varieties, and suppose that $H^\bullet(Y; \QQ)$ is spanned by classes of algebraic cycles, for instance $Y$ cellular (Example~\ref{hodge-conjecture:example-cellular}). Then
$\mathrm{HC}(X \times Y)$ holds if and only if $\mathrm{HC}(X)$ holds.
\end{proposition}

\begin{proof}
Choose a basis $y_j = \mathrm{cl}(W_j) \in H^{2k_j}(Y; \QQ)$ of algebraic classes; they are of type $(k_j, k_j)$ (Corollary~\ref{intersection-theory:corollary-hodge-classes}), and $H^\bullet(Y; \QQ)$ has no odd part. By the
K\"unneth formula (Theorem~\ref{singular-homology:theorem-kunneth}), every $\beta \in H^{2p}(X \times Y; \QQ)$ is uniquely $\beta = \sum_j\mathrm{pr}_1^*\alpha_j \cup \mathrm{pr}_2^*y_j$ with
$\alpha_j \in H^{2p-2k_j}(X; \QQ)$, and the same holds over $\CC$. Since pullbacks and cup products preserve types (Proposition~\ref{kahler:proposition-functoriality}), the component of type $(r, s)$ of
$\alpha_j$ contributes a class of type $(r + k_j, s + k_j)$. If $\beta$ is of type $(p, p)$, comparing components of type $(p, p)$ and using the uniqueness of the decomposition shows that each $\alpha_j$ is of type
$(p - k_j, p - k_j)$, hence a Hodge class. If $\mathrm{HC}(X)$ holds, $\alpha_j = \mathrm{cl}(Z_j)$ for a rational cycle $Z_j$, and $\beta = \sum_j\mathrm{cl}(Z_j \times W_j)$ is algebraic
(Theorem~\ref{intersection-theory:theorem-cycle-class}(1) and (2)).

Conversely, let $\alpha \in \mathrm{Hdg}^p(X)$ and let $y_0 = \mathrm{cl}(\mathrm{pt}) \in H^{2\dim Y}(Y; \QQ)$. The class $\mathrm{pr}_1^*\alpha \cup \mathrm{pr}_2^*y_0$ is a Hodge class on $X \times Y$; it is the pullback of $\alpha$ to $X \times \{\mathrm{pt}\}$
pushed forward to $X \times Y$, and if it is algebraic, so is its pushforward $\mathrm{pr}_{1*}(\mathrm{pr}_1^*\alpha \cup \mathrm{pr}_2^*y_0) = \alpha \cup \mathrm{pr}_{1*}\mathrm{pr}_2^*y_0 = \alpha$, by the projection formula and
$\mathrm{pr}_{1*}\mathrm{pr}_2^*y_0 = 1$.
\end{proof}

\begin{example}\label{known-cases:example-cellular-factor}
The Hodge conjecture holds for $S \times \PP^m$, $S \times \mathrm{Gr}(k, n)$ and $S \times F$ with $F$ a flag variety, for every nonsingular projective surface or threefold $S$, although these products can have
arbitrary dimension: by Proposition~\ref{known-cases:proposition-cellular-factor} and Theorem~\ref{known-cases:theorem-general}(1),(2). The same argument shows that the Hodge classes of $X \times Y$ are all products of Hodge
classes of $X$ with classes of $Y$ in this situation; the interesting Hodge classes on products come from factors with odd cohomology, as for $E \times E'$
(Example~\ref{hodge-conjecture:example-elliptic-product}).
\end{example}

\section{Abelian varieties}\label{known-cases:section-abelian}

\begin{remark}\label{known-cases:remark-abelian-mechanism}
For an abelian variety $A$, the Hodge classes on all powers $A^r$ are the invariants of the Hodge group $\mathrm{Hg} = \mathrm{Hg}(H^1(A))$ (Corollary~\ref{hodge-structures:corollary-hodge-ring-abelian}),
and the divisor classes on $A^r$ are the invariants of the \emph{Lefschetz group} $\mathrm{Lef}(A)$, the connected centraliser of $\End(A) \otimes \QQ$ in the symplectic group of a polarisation
\cite{murty-1984}. Always $\mathrm{Hg} \subseteq \mathrm{Lef}(A)$. If equality holds, the Hodge ring of every power of $A$ is generated by divisors, and the Hodge conjecture holds for all powers of $A$.
\end{remark}

\begin{theorem}\label{known-cases:theorem-abelian-known}
$\mathrm{Hg}(A) = \mathrm{Lef}(A)$, and hence the Hodge conjecture holds for all powers of $A$, in each of the following cases: $A$ is a product of elliptic curves \cite{imai-1976},
\cite{murty-1984}; $A$ is simple of prime dimension (Tankeev) \cite{tankeev-1982}; $A$ is of dimension at most $5$ and has no factor of Weil type (Moonen--Zarhin) \cite{moonen-zarhin-1999}. More
generally the Hodge conjecture holds for $A$ itself (not necessarily its powers) whenever $\mathrm{Hdg}^\bullet(A) = \mathrm{Lef}^\bullet(A)$, which holds for $A$ of dimension $\leq 3$ by
Theorem~\ref{known-cases:theorem-general}(1).
\end{theorem}

\begin{proposition}[Weil classes]\label{known-cases:proposition-weil-classes}
Let $A$ be a complex abelian fourfold and $K \subseteq \End(A) \otimes \QQ$ an imaginary quadratic field such that, for the two embeddings $\sigma, \bar\sigma \colon K \to \CC$, the eigenspaces of $K$
on $H^{1,0}(A)$ both have dimension $2$ (\emph{Weil type}). Then $H^1(A; \QQ)$ is a $4$-dimensional $K$-vector space. Let $V_\sigma, V_{\bar\sigma} \subseteq H^1(A; \CC)$ be the eigenspaces of $K$.
There is a $2$-dimensional $\QQ$-subspace $W_K \subseteq \Lambda^4H^1(A; \QQ) = H^4(A; \QQ)$ with $W_K \otimes \CC = \Lambda^4V_\sigma \oplus \Lambda^4V_{\bar\sigma}$, and it consists of Hodge classes, the
\emph{Weil classes}. For a general abelian fourfold of Weil type, $\mathrm{Hdg}^2(A) = \mathrm{Lef}^2(A) \oplus W_K$ with $W_K \cap \mathrm{Lef}^2(A) = 0$, so the Weil classes are
not combinations of intersections of divisors.
\end{proposition}

\begin{proof}
$K$ acts on $H_\CC = H^1(A; \CC)$ with eigenspaces $V_\sigma$, $V_{\bar\sigma}$, each of dimension $4$, and $K$ acts by Hodge endomorphisms, so $V_\sigma = V_\sigma^{1,0} \oplus V_\sigma^{0,1}$. By
hypothesis $\dim V_\sigma^{1,0} = 2$, and then $\dim V_\sigma^{0,1} = 2$. The space $\Lambda^4_KH^1(A; \QQ) \otimes_\QQ\CC$ decomposes as $\Lambda^4V_\sigma \oplus \Lambda^4V_{\bar\sigma}$, and
$\Lambda^4V_\sigma = \Lambda^2V_\sigma^{1,0} \otimes \Lambda^2V_\sigma^{0,1}$ is of type $(2,2)$, and likewise for $\bar\sigma$. The subspace $\Lambda^4V_\sigma \oplus \Lambda^4V_{\bar\sigma}$ of
$\Lambda^4H^1(A; \CC)$ is defined over $K$, since the eigenspace decomposition is, and it is stable under the nontrivial automorphism of $K$, which exchanges $V_\sigma$ and $V_{\bar\sigma}$. By Galois
descent it is the complexification of a $\QQ$-subspace $W_K$, of dimension $2$, whose elements are rational classes of type $(2,2)$. That $W_K \cap \mathrm{Lef}^2 = 0$ for general $A$ of Weil type, where $\End(A) \otimes \QQ = K$ and the Lefschetz
group is a unitary group $\mathrm{U}(2,2)$ while $\mathrm{Hg} = \mathrm{SU}(2,2)$, is \cite{weil-1977}; see also \cite{birkenhake-lange}.
\end{proof}

\begin{remark}\label{known-cases:remark-weil-classes}
Weil classes are the essential difficulty for abelian varieties: they are the first Hodge classes not generated by divisors, and Moonen and Zarhin showed that in dimension $\leq 5$
all exceptional Hodge classes come from Weil classes. The Hodge conjecture for Weil classes was proved by Schoen for $K = \QQ(\sqrt{-3})$ and $K = \QQ(i)$ in special cases \cite{schoen-1988},
by van Geemen in further cases, and by Markman for all abelian fourfolds of Weil type of discriminant $1$ \cite{markman-2023}. Weil classes are absolute Hodge classes by Deligne's theorem
(Chapter~\ref{absolute-hodge}).
\end{remark}

\section{Hypersurfaces and complete intersections}\label{known-cases:section-hypersurfaces}

\begin{proposition}\label{known-cases:proposition-hypersurfaces}
Let $X \subseteq \PP^{n+1}$ be a nonsingular hypersurface, or more generally a nonsingular complete intersection of dimension $n$ in some $\PP^N$, and $h$ the hyperplane class. Then for $2p \neq n$,
$H^{2p}(X; \QQ) = \QQ h^p$. Hence $\mathrm{HC}(X, p)$ holds for $2p \neq n$.
\end{proposition}

\begin{proof}
By the Lefschetz hyperplane theorem (Theorem~\ref{kodaira:theorem-lefschetz-hyperplane}), applied repeatedly for complete intersections, $H^k(X; \QQ) \cong H^k(\PP^N; \QQ)$ for $k < n$, which is
$\QQ h^{k/2}$ for $k$ even. For $k > n$, Poincar\'e duality gives $b_k = b_{2n-k}$, which is $1$ for $k$ even, and $h^{k/2} \neq 0$ by Proposition~\ref{kahler:proposition-kahler-class-powers}, so
$H^k = \QQ h^{k/2}$. Powers of $h$ are algebraic, being classes of complete intersections of hyperplane sections.
\end{proof}

\begin{remark}\label{known-cases:remark-middle-degree}
In the middle degree $n = 2p$, a very general hypersurface of degree $d$ has $\mathrm{Hdg}^p(X) = \QQ h^p$ as soon as $H^{n}$ is not of Tate type up to $h^p$, for instance for surfaces with $d \geq 4$
(Theorem~\ref{hodge-conjecture:theorem-noether-lefschetz}) and for fourfolds with $d \geq 3$, so the conjecture holds there trivially. Special hypersurfaces are the interesting cases. For cubic
fourfolds, Zucker proved the conjecture with normal functions (Remark~\ref{vhs:remark-poincare-lefschetz}), and Voisin proved the integral version \cite{voisin-2006}. For
Fermat hypersurfaces $x_0^d + \cdots + x_{n+1}^d = 0$, the Hodge classes are computable via the action of $(\mu_d)^{n+2}$ and characters (Jacobi sums), and Shioda proved the conjecture for many $(d, n)$,
including $d$ prime and $d = 4$ \cite{shioda-1979}.
\end{remark}

\section{Uniruled varieties and the decomposition of the diagonal}\label{known-cases:section-diagonal}

\begin{theorem}[Bloch--Srinivas]\label{known-cases:theorem-bloch-srinivas}
Suppose that $\CH_0(X)_\QQ$ is supported on a closed subvariety of dimension $\leq 3$, that is, there is $W \subseteq X$ of dimension $\leq 3$ with $\CH_0(W)_\QQ \to \CH_0(X)_\QQ$ surjective. Then
$\mathrm{HC}(X, 2)$ holds. This applies to rationally connected varieties, for which $\CH_0(X)_\QQ = \QQ$.
\end{theorem}

This is \cite{bloch-srinivas-1983}. The method is the \emph{decomposition of the diagonal}: the hypothesis gives an equality $N\Delta_X = Z_1 + Z_2$ in $\CH^n(X \times X)$ with $Z_1$ supported on
$D \times X$ for a divisor $D$ and $Z_2$ supported on $X \times W$. Letting both sides act on $H^4(X; \QQ)$ as correspondences, the identity (times $N$) factors through the cohomology of a resolution of
$D$, where Hodge classes of degree $2$ on a variety of lower dimension are handled by induction and the $(1,1)$-theorem, plus a contribution through $W$, of dimension $\leq 3$, where the conjecture is
known. The same argument gives the conjecture for uniruled fourfolds, first proved by Conte and Murre \cite{conte-murre-1978}.

\section{K3 surfaces, hyper-K\"ahler varieties and Kuga--Satake}\label{known-cases:section-k3}

\begin{remark}\label{known-cases:remark-k3}
For a single K3 surface the conjecture holds, by dimension. The interesting questions are on products. \emph{Mukai's conjecture}, proved by Buskin \cite{buskin-2019}, says that every rational Hodge
isometry $H^2(S; \QQ) \to H^2(S'; \QQ)$ between K3 surfaces, a Hodge class in $H^4(S \times S'; \QQ)$, is algebraic. Markman extended this to hyper-K\"ahler varieties of $K3^{[n]}$-type.
\end{remark}

\begin{remark}[Kuga--Satake]\label{known-cases:remark-kuga-satake}
Kuga and Satake attached to a polarised K3 surface $S$ an abelian variety $\mathrm{KS}(S)$, of dimension $2^{19}$, built from the even Clifford algebra of the lattice of primitive classes, together with
an embedding of Hodge structures $H^2_{\mathrm{prim}}(S; \QQ) \hookrightarrow H^1(\mathrm{KS}(S)) \otimes H^1(\mathrm{KS}(S))$ \cite{kuga-satake-1967}. The Hodge conjecture predicts that this embedding is
induced by an algebraic correspondence between $S$ and $\mathrm{KS}(S)^2$ (the \emph{Kuga--Satake Hodge conjecture}). It is known for some families, for instance by Paranjape for double planes branched
along six lines \cite{paranjape-1988}. Deligne used the Kuga--Satake construction, which works in families and in characteristic $p$, to prove the Weil conjectures for K3 surfaces
\cite{deligne-k3}, and to show that Hodge classes on K3 surfaces and their products with abelian varieties behave like absolute Hodge classes.
\end{remark}

\section{The evidence assessed}\label{known-cases:section-assessment}

\begin{remark}\label{known-cases:remark-assessment}
The known cases share a feature: the Hodge classes are either generated by divisors, or produced by an explicit geometric construction adapted to the variety (normal functions for cubic fourfolds,
the decomposition of the diagonal, families of cycles on special abelian fourfolds, derived equivalences of K3 surfaces). There is no general mechanism. At the same time no case has failed, the
Hodge loci behave as the conjecture predicts (Theorem~\ref{vhs:theorem-cdk}), and Hodge classes on abelian varieties satisfy the arithmetic consequence of the conjecture (Deligne,
Chapter~\ref{absolute-hodge}). The first genuinely open cases include Weil classes on general abelian fourfolds of Weil type of arbitrary discriminant, abelian varieties of dimension $\geq 6$ with
exceptional classes, and products of two K3 surfaces in degree $4$ away from Hodge isometries.
\end{remark}

\section{Exercises}\label{known-cases:section-exercises}

\begin{exercise}\label{known-cases:exercise-cubic-fourfold-numbers}
For a nonsingular cubic fourfold $X \subseteq \PP^5$, show that $H^4(X; \QQ)$ has Hodge numbers $(0, 1, 21, 1, 0)$ and that $\mathrm{HC}(X, p)$ holds for $p \neq 2$ by
Proposition~\ref{known-cases:proposition-hypersurfaces}.
\end{exercise}

\begin{solution}
As in Proposition~\ref{calabi-yau:proposition-quintic}, $c(TX) = (1 + h)^6/(1 + 3h)$, whose degree-$4$ part is $(15 - 3 \cdot 20 + 9 \cdot 15 - 27 \cdot 6 + 81)h^4 = 9h^4$, and
$\int_Xh^4 = 3$, so $\chi(X) = 27$. The odd Betti numbers vanish and $b_0 = b_2 = b_6 = b_8 = 1$ (Proposition~\ref{known-cases:proposition-hypersurfaces}), so $b_4 = 27 - 4 = 23$. The Hodge
numbers $h^{4,0} = 0$, $h^{3,1} = 1$, $h^{2,2} = 21$ follow from Griffiths' description of the primitive cohomology of hypersurfaces by residues \cite{voisin-hodge-2}. The second statement is
Proposition~\ref{known-cases:proposition-hypersurfaces} with $n = 4$.
\end{solution}

\begin{exercise}\label{known-cases:exercise-lefschetz-ring-product}
Let $E$ be an elliptic curve without complex multiplication. Show that $\mathrm{Hdg}^\bullet(E^r) = \mathrm{Lef}^\bullet(E^r)$ for all $r$.
\end{exercise}

\begin{solution}
By Proposition~\ref{hodge-structures:proposition-elliptic-curves}, $\mathrm{Hg}(E) = \mathrm{SL}_2$, and the invariants of $\mathrm{SL}_2$ acting diagonally on $\Lambda^\bullet(H^{\oplus r})$ are generated, by classical
invariant theory, by the invariants in $\Lambda^2(H^{\oplus r})$, which are the divisor classes (Exercise~\ref{algebraic-groups:exercise-hodge-classes-torus}). Apply
Corollary~\ref{hodge-structures:corollary-hodge-ring-abelian}.
\end{solution}

\section{Variants That Fail}\label{known-cases:section-failures}
The Hodge conjecture is surrounded by statements that are equally natural and demonstrably false. Each failure identifies a mechanism that Hodge theory does not
control, and so a constraint on any proof. We treat four variants: the integral Hodge conjecture, false for torsion classes (Atiyah--Hirzebruch) and for non-torsion classes
(Koll\'ar); the K\"ahler version, false even with Chern classes of coherent sheaves (Voisin); Hodge's general conjecture, false for formal reasons and repaired by Grothendieck;
and the variational Hodge conjecture, which is open but belongs here as the statement one would need to propagate algebraicity in families. The chapter ends with a diagnostic list.
References: the source \emph{The Hodge Conjecture} (Chapter 7), \cite{atiyah-hirzebruch-1962}, \cite{kollar-1992}, \cite{voisin-2002}, \cite{grothendieck-1969}.

\noindent\textbf{Prerequisites.} Chapter~\ref{known-cases} (Known Cases and Evidence).

Throughout, $X$ is a nonsingular complex projective variety of dimension $n$ unless stated otherwise.

\section{The integral Hodge conjecture}\label{known-cases:section-integral}

\begin{definition}\label{known-cases:definition-defect}
The \emph{integral Hodge conjecture} $\mathrm{IHC}(X, p)$ asserts that the group $\mathrm{Hdg}^p(X; \ZZ)$ of integral Hodge classes (Definition~\ref{intersection-theory:definition-hodge-classes}) is generated by
classes of subvarieties. Its \emph{defect} is the group $Z^{2p}(X) = \mathrm{Hdg}^p(X; \ZZ)/\mathrm{cl}(\CH^p(X))$.
\end{definition}

\begin{proposition}\label{known-cases:proposition-defect-finite}
$Z^{2p}(X)$ is finitely generated, and it is finite if and only if $\mathrm{HC}(X, p)$ holds. $Z^2(X) = 0$, and $Z^{2n}(X) = 0$.
\end{proposition}

\begin{proof}
$\mathrm{Hdg}^p(X; \ZZ)$ is a subgroup of the finitely generated group $H^{2p}(X; \ZZ)$. The group $Z^{2p}(X) \otimes \QQ$ is $\mathrm{Hdg}^p(X)/Z^p_{\mathrm{alg}}(X)$, which vanishes exactly when
$\mathrm{HC}(X, p)$ holds, and a finitely generated group is finite if and only if its rationalisation vanishes. $Z^2 = 0$ is the Lefschetz $(1,1)$-theorem (Theorem~\ref{hodge-conjecture:theorem-lefschetz-11}),
and $H^{2n}(X; \ZZ)$ is generated by the classes of points, one for each component.
\end{proof}

\begin{theorem}[Atiyah--Hirzebruch]\label{known-cases:theorem-atiyah-hirzebruch}
There are nonsingular projective varieties $X$ with torsion classes in $H^4(X; \ZZ)$, automatically integral Hodge classes, that are not algebraic. In particular $\mathrm{IHC}$ fails.
\end{theorem}

This is \cite{atiyah-hirzebruch-1962}. The obstruction is topological. A subvariety $Z$ of codimension $p$ defines an element $[\mathcal{O}_Z]$ of topological $K$-theory $K^0(X)$, through a finite
resolution by vector bundles, and its image in the graded piece of the topological filtration is $\mathrm{cl}(Z)$ up to the factor $(p-1)!$, since $c_p(\mathcal{O}_Z) = (-1)^{p-1}(p-1)!\,\mathrm{cl}(Z)$
while $\mathrm{ch}_p(\mathcal{O}_Z) = \mathrm{cl}(Z)$ (Proposition~\ref{intersection-theory:proposition-chern-character-cycle}). So algebraic classes survive to $E_\infty$ in the Atiyah--Hirzebruch spectral sequence
$H^s(X; K^t(\mathrm{pt})) \Rightarrow K^{s+t}(X)$, whose first possibly nonzero differential is $d_3 = \mathrm{Sq}^3_\ZZ = \beta\circ\mathrm{Sq}^2\circ\rho$, the integral lift of the Steenrod square
(Chapter~\ref{cohomology-operations}). Hence an algebraic class $\alpha$ satisfies $\mathrm{Sq}^3_\ZZ(\alpha) = 0$. Atiyah and Hirzebruch then produce, by the Godeaux--Serre construction of a quotient
$X = Y/G$ of a complete intersection by a free action of a finite group $G$ such as $(\ZZ/p)^3$, whose cohomology agrees with that of the classifying space $BG$ in low degrees, a torsion class with
$\mathrm{Sq}^3_\ZZ \neq 0$.

\begin{remark}\label{known-cases:remark-totaro}
Totaro observed that the cycle class map factors through complex cobordism, $\CH^p(X) \to MU^{2p}(X) \otimes_{MU^*}\ZZ \to H^{2p}(X; \ZZ)$, which gives a finer obstruction
\cite{totaro-1997}. Soul\'e and Voisin produced further torsion counterexamples by other methods \cite{soule-voisin-2005}. All these classes vanish in $H^{2p}(X; \QQ)$, so they do not affect the rational
conjecture: Hodge theory is blind to torsion.
\end{remark}

\begin{theorem}[Koll\'ar]\label{known-cases:theorem-kollar}
Let $\ell \geq 5$ be a prime and $d$ a multiple of $\ell^3$. For a very general hypersurface $X \subseteq \PP^4$ of degree $d$, every curve $C \subseteq X$ has degree divisible by $\ell$. Consequently the
generator $\alpha$ of $H^4(X; \ZZ) \cong \ZZ$ is an integral Hodge class that is not algebraic, and $Z^4(X) \neq 0$ contains an element of order $\ell$.
\end{theorem}

\begin{proof}[Proof of the consequence]
The divisibility statement is \cite{kollar-1992}. Granting it: $\dim X = 3$, and by the Lefschetz hyperplane theorem (Theorem~\ref{kodaira:theorem-lefschetz-hyperplane}) and Poincar\'e duality,
$H^2(X; \ZZ) \cong \ZZ h$ and $H^4(X; \ZZ) \cong \ZZ$, generated by a class $\alpha$ with $\int_X\alpha \cup h = 1$ (the pairing $H^2 \times H^4 \to \ZZ$ is perfect on the torsion-free quotients; integral
Lefschetz for $H^2$ uses the integral version of the hyperplane theorem, Remark~\ref{kodaira:remark-lefschetz-integral}). Since $H^4(X; \CC) = \CC h^2$ is of type $(2,2)$, $\alpha$ is an integral Hodge class. A
$1$-cycle $\sum a_iC_i$ has class $(\sum a_i\deg C_i)\alpha$, because $\int\mathrm{cl}(C) \cup h = \deg C$ (Theorem~\ref{intersection-theory:theorem-cycle-class}(4)). If all degrees are divisible by $\ell$, the algebraic
classes lie in $\ell\ZZ\alpha$, so $\alpha$ is not algebraic and $\ZZ/\ell$ is a quotient of $Z^4(X)$.
\end{proof}

\begin{corollary}\label{known-cases:corollary-kollar-defect}
In the situation of Theorem~\ref{known-cases:theorem-kollar}, the defect $Z^4(X)$ is a finite cyclic group whose order is divisible by $\ell$ and divides $d$. In particular $\mathrm{HC}(X, 2)$ holds although
$\mathrm{IHC}(X, 2)$ fails.
\end{corollary}

\begin{proof}
With $\alpha$ the generator of $H^4(X; \ZZ) = \ZZ\alpha$ normalised by $\int_X\alpha \cup h = 1$, the class $h^2$ satisfies $\int_Xh^2 \cup h = \deg X = d$, so $h^2 = d\alpha$. It is algebraic, the class of the intersection of $X$
with two hyperplanes, a curve of degree $d$. So $d\ZZ\alpha \subseteq \mathrm{cl}(\CH^2(X)) \subseteq \ell\ZZ\alpha$, the second inclusion by Theorem~\ref{known-cases:theorem-kollar}. Hence
$\mathrm{cl}(\CH^2(X)) = m\ZZ\alpha$ with $\ell \mid m \mid d$, and $Z^4(X) = \mathrm{Hdg}^2(X; \ZZ)/\mathrm{cl}(\CH^2(X)) = \ZZ\alpha/m\ZZ\alpha \cong \ZZ/m$. Rationally $\alpha = \frac1dh^2$ is algebraic.
\end{proof}

\begin{remark}\label{known-cases:remark-kollar-mechanism}
The divisibility is proved by specialisation and countability. The relative Hilbert scheme of curves in the universal family of hypersurfaces has countably many components, each proper over the
parameter space. For a very general member $X$, every curve on $X$ therefore lies in a component that dominates the parameter space, so it specialises to a $1$-cycle of the same degree on every
member, in particular on a special singular member $X_0$. Koll\'ar builds $X_0$ from a cyclic covering of degree $\ell$ so that every $1$-cycle on it has degree divisible by $\ell$; this is where
$\ell^3 \mid d$ is used. (Koll\'ar's statement allows any integer $\ell$ prime to $6$.) The Hodge structure here is as simple as possible; the obstruction is purely about degrees of curves, on which Hodge theory has nothing to say.
\end{remark}

\begin{theorem}\label{known-cases:theorem-what-survives}
\begin{enumerate}
\item (Voisin) $\mathrm{IHC}(X, 2)$ holds for uniruled threefolds and for Calabi--Yau threefolds \cite{voisin-2006}.
\item (Colliot-Th\'el\`ene--Voisin) The group $Z^4(X)$ is a birational invariant of nonsingular projective varieties. Its torsion subgroup is a quotient of the unramified cohomology group
$H^3_{\mathrm{nr}}(X, \QQ/\ZZ(2))$, and is isomorphic to it when $\CH_0(X)$ is supported on a surface, for instance when $X$ is rationally connected. Since unramified cohomology vanishes for
rational varieties, a nonzero torsion element of $Z^4(X)$ obstructs rationality \cite{colliot-thelene-voisin-2012}.
\end{enumerate}
\end{theorem}

\section{The K\"ahler counterexample}\label{known-cases:section-kahler}

\begin{theorem}[Voisin]\label{known-cases:theorem-voisin-kahler}
There is a complex torus $T$ of dimension $4$ with a nonzero class in $\mathrm{Hdg}^2(T)$ that is not in the $\QQ$-span of the classes $\mathrm{ch}_2(\mathcal{F})$ of coherent analytic sheaves on $T$, nor of classes of
analytic subvarieties. Hence the Hodge conjecture fails for compact K\"ahler manifolds in both the cycle form and the sheaf form.
\end{theorem}

This is \cite{voisin-2002}. The torus is of Weil type (Proposition~\ref{known-cases:proposition-weil-classes}): the construction of Weil classes uses only Hodge theory, not projectivity. In outline,
Voisin chooses the lattice so that $\mathrm{NS}(T) = 0$ and the Hodge classes in $H^4$ are primitive for a K\"ahler class $\omega$, that is, $\int_T\alpha \wedge \omega^2 = 0$. For a coherent sheaf
$\mathcal{F}$ on $T$, $c_1(\mathcal{F}) = 0$ because $\mathrm{NS}(T) = 0$, and a positivity property of $c_2$ coming from Hermite--Einstein metrics (a Bogomolov-type inequality) gives
$\int_Tc_2(\mathcal{F}) \wedge \omega^2 \geq 0$ with equality only if $c_2(\mathcal{F}) = 0$. Since $c_2(\mathcal{F})$ is a Hodge class, it is primitive, so $c_2(\mathcal{F}) = 0$. Analytic subvarieties are
covered because their structure sheaves are coherent sheaves.

\begin{remark}\label{known-cases:remark-kahler-lesson}
All the Hodge-theoretic structure (Hodge decomposition, hard Lefschetz, Hodge--Riemann relations, the $\partial\bar\partial$-lemma) holds on Voisin's torus. So no proof of the Hodge conjecture can use
only the K\"ahler package: it must use projectivity, through sections of ample line bundles, algebraic Hilbert schemes, GAGA, or arithmetic models. For $p = 1$ the K\"ahler version holds in the form
of Theorem~\ref{hodge-conjecture:theorem-kahler}: integral $(1,1)$-classes are Chern classes of holomorphic line bundles.
\end{remark}

\section{Hodge's general conjecture}\label{known-cases:section-general-failures}

\begin{proposition}\label{known-cases:proposition-general-false}
Hodge's original general conjecture $N^cH^k(X; \QQ) = H^k(X; \QQ) \cap F^cH^k(X; \CC)$ fails in general: the left side is a sub-Hodge structure
(Proposition~\ref{hodge-conjecture:proposition-coniveau}), while the right side need not be.
\end{proposition}

\begin{proof}
Let $V = H^k(X; \QQ)$ and $W = V \cap F^c$. Since $V$ is stable under conjugation, $W \subseteq F^c \cap \overline{F^c}$, which is $\bigoplus_{c \leq p \leq k-c}V^{p,k-p}$. A rational subspace of this sum need
not be the rational form of a sum of Hodge summands, that is, need not be a sub-Hodge structure. Grothendieck exhibits an abelian variety whose $H^3$ contains a rational line in $F^1$ whose smallest
containing sub-Hodge structure is not contained in $F^1$ \cite{grothendieck-1969}. For such a line, the right side contains it, and the left side, a sub-Hodge structure contained in $F^1$
(Proposition~\ref{hodge-conjecture:proposition-coniveau}), cannot.
\end{proof}

\begin{remark}\label{known-cases:remark-general-lesson}
The lesson is formal: the objects of Hodge theory are sub-Hodge structures, not rational subspaces satisfying a filtration condition. Grothendieck's repaired \emph{generalised Hodge conjecture}
(Remark~\ref{hodge-conjecture:remark-grothendieck}) is open in essentially every case where the Hodge conjecture is, and in many where it is known.
\end{remark}

\section{The variational Hodge conjecture}\label{known-cases:section-variational}

\begin{conjecture}[Variational Hodge conjecture]\label{known-cases:conjecture-variational}
Let $f \colon \mathcal{X} \to S$ be a smooth projective morphism over a connected nonsingular base, and $\lambda$ a flat section of $R^{2p}f_*\QQ$ with $\lambda_s \in \mathrm{Hdg}^p(X_s)$ for all $s$. If $\lambda_{s_0}$ is
algebraic for one $s_0$, then $\lambda_s$ is algebraic for all $s$.
\end{conjecture}

\begin{remark}\label{known-cases:remark-variational}
The variational conjecture follows from the Hodge conjecture applied to a compactification of the total space, by the theorem of the fixed part (Theorem~\ref{vhs:theorem-fixed-part}). Bloch's
semiregularity theorem proves the infinitesimal version: a local complete intersection cycle whose semiregularity map is injective deforms together with its cohomology class
\cite{bloch-1972}. The Cattani--Deligne--Kaplan theorem (Theorem~\ref{vhs:theorem-cdk}) says that the locus where $\lambda$ stays Hodge is algebraic; the variational conjecture says that on this locus
$\lambda$ stays algebraic. The first is a theorem and the second is open, which summarises the state of the subject.
\end{remark}

\section{A diagnostic list}\label{known-cases:section-diagnostic}

\begin{remark}\label{known-cases:remark-diagnostic}
A proof of the Hodge conjecture must be able to distinguish it from its false neighbours:
\begin{enumerate}
\item \emph{It must lose denominators.} Since the integral version is false (Koll\'ar), cycles can only be produced after clearing denominators. In known cases the denominators enter through
push-pull, $f_*f^* = \deg f$, or through factorials such as $c_p(\mathcal{O}_Z) = (-1)^{p-1}(p-1)!\,\mathrm{cl}(Z)$.
\item \emph{It must use projectivity}, not just the K\"ahler package (Voisin).
\item \emph{It must be insensitive to torsion} (Atiyah--Hirzebruch): it cannot produce integral cycles from integral Hodge classes.
\item \emph{It must respect sub-Hodge structures} (Grothendieck): it must work with Hodge classes, not with rational classes in a filtration step.
\end{enumerate}
\end{remark}

\section{Exercises on Variants That Fail}\label{known-cases:section-exercises-failures}

\begin{exercise}\label{known-cases:exercise-torsion-hodge}
Show that every torsion class in $H^{2p}(X; \ZZ)$ is an integral Hodge class, and that $Z^{2p}(X)$ is torsion if $\mathrm{HC}(X, p)$ holds.
\end{exercise}

\begin{solution}
A torsion class maps to $0 \in H^{2p}(X; \CC)$, which lies in $H^{p,p}$. The second statement is Proposition~\ref{known-cases:proposition-defect-finite}.
\end{solution}

\begin{exercise}\label{known-cases:exercise-kollar-quintic}
Why does Theorem~\ref{known-cases:theorem-kollar} not contradict the fact that a quintic threefold contains lines, which have degree $1$?
\end{exercise}

\begin{solution}
The theorem needs $d$ divisible by $\ell^3 \geq 125$, and the hypersurface very general. A quintic has $d = 5$, which is not divisible by $125$.
\end{solution}
